% PPMI-Pipeline v2.1.6 poster source.
% Compile with pdfLaTeX, LuaLaTeX, or XeLaTeX.
% Required local theme files:
%   beamerthemegemini.sty
%   beamercolorthemegemini.sty
%   beamerinnerthemegemini.sty
% Required poster figures and QR assets belong in figures/.

\begin{filecontents*}{poster.bib}
@misc{ppmi_dataset,
  author       = {{Parkinson's Progression Markers Initiative}},
  title        = {{Parkinson's Progression Markers Initiative (PPMI) database}},
  year         = {2026},
  howpublished = {\url{https://www.ppmi-info.org/access-data-specimens/download-data}},
  note         = {RRID:SCR\_006431}
}

@inproceedings{mckinney2010,
  author    = {McKinney, Wes},
  title     = {Data Structures for Statistical Computing in Python},
  booktitle = {Proceedings of the 9th Python in Science Conference},
  pages     = {56--61},
  year      = {2010},
  doi       = {10.25080/Majora-92bf1922-00a}
}

@article{pedregosa2011,
  author  = {Pedregosa, Fabian and Varoquaux, Ga\"{e}l and Gramfort, Alexandre
             and Michel, Vincent and Thirion, Bertrand and Grisel, Olivier
             and Blondel, Mathieu and Prettenhofer, Peter and Weiss, Ron
             and Dubourg, Vincent and Vanderplas, Jake and Passos, Alexandre
             and Cournapeau, David and Brucher, Matthieu and Perrot, Matthieu
             and Duchesnay, {\'{E}douard}},
  title   = {Scikit-learn: Machine Learning in Python},
  journal = {Journal of Machine Learning Research},
  volume  = {12},
  number  = {85},
  pages   = {2825--2830},
  year    = {2011}
}
\end{filecontents*}

\documentclass[final]{beamer}

\usepackage[T1]{fontenc}
\usepackage[utf8]{inputenc}
\usepackage{lmodern}
\usepackage[size=custom,width=120,height=72,scale=1.0]{beamerposter}
\usetheme{gemini}
\usecolortheme{gemini}

\usepackage{graphicx}
\usepackage{booktabs}
\usepackage[numbers]{natbib}
\usepackage{url}
\usepackage{tikz}
\usepackage{pgfplots}
\usepackage{anyfontsize}

\pgfplotsset{compat=1.14}
\graphicspath{{figures/}{./}}
\pdfstringdefDisableCommands{%
  \def\translate#1{#1}%
}

% Print-scale typography for a 120 cm x 72 cm conference poster.
\setbeamerfont{title in headline}{series=\bfseries,size=\fontsize{50}{56}}
\setbeamerfont{author in headline}{size=\fontsize{27}{32}}
\setbeamerfont{institute in headline}{size=\fontsize{21}{26}}
\setbeamerfont{block title}{series=\bfseries,size=\fontsize{27}{32}}
\setbeamerfont{block body}{size=\fontsize{21}{27}}
\setbeamerfont{caption}{size=\fontsize{15}{19}}
\setbeamercolor{metrics box}{fg=white,bg=geminiTeal}
\setbeamerfont{metrics box}{series=\bfseries,size=\fontsize{20}{25}}

% Three-column layout: (N + 1) * sepwidth + N * colwidth = paperwidth.
\newlength{\sepwidth}
\newlength{\colwidth}
\setlength{\sepwidth}{0.025\paperwidth}
\setlength{\colwidth}{0.3\paperwidth}
\newcommand{\separatorcolumn}{\begin{column}{\sepwidth}\end{column}}

\title{Reproducible Fixed-Horizon Prognostic Modelling of Parkinson's Disease Using Longitudinal PPMI Data}
\author{Faris Haq}
\institute{Data Science \& Biomedical Research}

\begin{document}

\begin{frame}[t]
\begin{columns}[t]

\separatorcolumn

% -----------------------------------------------------------------------------
% Column 1: Objective and Methods
% -----------------------------------------------------------------------------
\begin{column}{\colwidth}

\begin{block}{Objective}
\begin{itemize}
  \setlength{\itemsep}{0.35em}
  \item \textbf{Clinical gap:} Longitudinal Parkinson's datasets are often analysed through opaque, one-off workflows with unclear cohort selection and temporal definitions.
  \item \textbf{Methodological risk:} Row-level evaluation can place observations from the same patient in both training and testing data, producing information leakage.
  \item \textbf{Aim:} Develop an auditable, reproducible fixed-horizon pipeline for prognostic modelling using PPMI clinical data \citep{ppmi_dataset}.
  \item \textbf{Scope:} Support transparent methodological research and patient stratification studies, not clinical diagnosis or treatment decisions.
\end{itemize}
\end{block}

\begin{block}{Methods}
\begin{itemize}
  \setlength{\itemsep}{0.35em}
  \item \textbf{Modular workflow:} \texttt{etl.py} validates and normalises PPMI exports; schema harmonisation standardises patient, visit, score, and date fields.
  \item \textbf{Fixed-horizon design:} The earliest usable baseline (BL) is identified for each patient, then the eligible follow-up nearest the prespecified horizon is selected.
  \item \textbf{Outcome construction:} Patient-level baseline score, target score, and score-change records are created; follow-up exclusions are explicitly recorded.
  \item \textbf{Modelling:} ElasticNet-regularised logistic regression implemented with scikit-learn \citep{pedregosa2011} predicts a configured binary progression label. The current baseline benchmark uses \texttt{Baseline\_Score} as its sole predictor.
  \item \textbf{Configuration:} Horizon, tolerance window, score endpoint, progression threshold, and model settings are version-controlled.
\end{itemize}
\end{block}

\begin{block}{PPMI Data Acknowledgement}
{\fontsize{13.5}{17}\selectfont
Data used in the preparation of this article was obtained on 2026-07-15 from
the Parkinson's Progression Markers Initiative (PPMI) database
(\href{https://www.ppmi-info.org/access-data-specimens/download-data}{%
\nolinkurl{www.ppmi-info.org/access-data-specimens/download-data}}),
RRID:SCR\_006431. For up-to-date information on the study, visit
\href{https://www.ppmi-info.org}{\nolinkurl{www.ppmi-info.org}}.

\vspace{0.35cm}

PPMI -- a public-private partnership -- is funded by the Michael J. Fox
Foundation for Parkinson's Research, and funding partners; including AbbVie,
Alamar Biosciences, Aligning Science Across Parkinson's (ASAP), Arrowhead
Pharma, Arvinas, AskBio, BIAL, BioArctic, Biohaven, BlueRock Therapeutics,
Bristol Myers Squibb, Calico Labs, Capsida Biotherapeutics, Critical Path
Institute, DaCapo Brainscience, Denali, Edmond J. Safra Foundation, Eli Lilly,
Gain Therapeutics, GE Healthcare, Genentech, GSK, Insitro, Johnson \& Johnson
Innovative Medicine, Lundbeck, Merck, Neumora, Neuron23, Novartis, Olink,
Regeneron, Roche, Sanofi, Tenvie, UCB, Vanqua Bio, Voyager Therapeutics, The
Weston Family Foundation.
}
\end{block}

\end{column}

\separatorcolumn

% -----------------------------------------------------------------------------
% Column 2: Results and empirical figures
% -----------------------------------------------------------------------------
\begin{column}{\colwidth}

\begin{block}{Results}
\begin{itemize}
  \setlength{\itemsep}{0.30em}
  \item Cohort derivation records validated observations, baseline eligibility, fixed-horizon retention, and exclusions at each filtering stage.
  \item Validation integrity and held-out performance are assessed across patient-isolated folds using AUC-ROC, precision, recall, and F1-score.
  \item \texttt{Baseline\_Score} was the sole predictor. Its coefficient remained positive across all five folds; this indicates single-coefficient consistency, not multi-feature importance.
  \item Low recall indicates limited sensitivity for progression detection at the current decision threshold.
  \item Results are interpreted as a methodological benchmark, not as a clinically actionable decision-support model.
\end{itemize}

\vspace{0.30cm}

\begin{beamercolorbox}[wd=\linewidth,rounded=true,shadow=false,sep=1.2ex,center]{metrics box}
  {\fontsize{24}{29}\selectfont
    \textbf{Retained cohort: $N = 2{,}875$}\quad
    ($62.3\%$ of $4{,}616$ usable baseline patients)
  }\\[0.45cm]
  \begin{tabular}{@{}cccc@{}}
    \textbf{AUC-ROC} & \textbf{Precision} & \textbf{Recall} & \textbf{F1-score} \\
    $0.695 \pm 0.023$ & $0.290 \pm 0.155$ & $0.034 \pm 0.016$ & $0.061 \pm 0.029$
  \end{tabular}
\end{beamercolorbox}

\vspace{0.30cm}

\begin{figure}
  \centering
  \includegraphics[width=0.74\linewidth]{cohort_flow_diagram.pdf}
  \caption{Cohort derivation from validated records to the fixed-horizon modelling cohort, including baseline eligibility, retained follow-up, and exclusions.}
  \label{fig:cohort-flow}
\end{figure}

\vspace{0.35cm}

\begin{figure}
  \centering
  \begin{minipage}[t]{0.48\linewidth}
    \centering
    \includegraphics[width=\linewidth,height=6.5cm,keepaspectratio]{validation_audit.pdf}\\[-0.1cm]
    {\scriptsize Patient-isolation audit}
  \end{minipage}
  \hfill
  \begin{minipage}[t]{0.48\linewidth}
    \centering
    \includegraphics[width=\linewidth,height=6.5cm,keepaspectratio]{fold_performance.pdf}\\[-0.1cm]
    {\scriptsize Held-out fold performance}
  \end{minipage}

  \vspace{0.25cm}

  \begin{minipage}[t]{0.48\linewidth}
    \centering
    \includegraphics[width=\linewidth,height=6.5cm,keepaspectratio]{feature_stability.pdf}\\[-0.1cm]
    {\scriptsize Feature stability}
  \end{minipage}
  \hfill
  \begin{minipage}[t]{0.48\linewidth}
    \centering
    \includegraphics[width=\linewidth,height=6.5cm,keepaspectratio]{coefficient_heatmap_clustered.pdf}\\[-0.1cm]
    {\scriptsize Coefficient direction across folds}
  \end{minipage}
  \caption{Validation integrity and model behaviour across patient-isolated folds. The audit confirms zero patient overlap; performance and coefficient panels show fold-level variation, directionality, and consistency.}
  \label{fig:validation-performance}
\end{figure}

\end{block}

\end{column}

\separatorcolumn

% -----------------------------------------------------------------------------
% Column 3: Validation, Conclusion, and References
% -----------------------------------------------------------------------------
\begin{column}{\colwidth}

\begin{block}{Validation}
\begin{itemize}
  \setlength{\itemsep}{0.35em}
  \item \textbf{Patient isolation:} GroupKFold groups observations by patient identifier, preventing repeated-patient leakage across folds.
  \item \textbf{Leakage audit:} Every fold is checked for train/test patient overlap; the real-data analysis confirmed zero overlapping patients.
  \item \textbf{Automated testing:} All 13 unit tests passed, covering schema harmonisation, fixed-horizon construction, target creation, modelling, validation, and visualisation.
  \item \textbf{Traceability:} Run provenance, cohort counts, fold metrics, coefficients, and generated figures are retained as reproducible artifacts.
\end{itemize}
\end{block}

\begin{block}{Conclusion}
\begin{itemize}
  \setlength{\itemsep}{0.35em}
  \item \textbf{Reproducibility:} The pipeline converts longitudinal PPMI preparation into a documented, version-controlled research protocol.
  \item \textbf{Methodological validity:} Explicit baseline definitions, fixed-horizon follow-up selection, and patient-isolated validation reduce avoidable analytical bias.
  \item \textbf{Clinical relevance:} Transparent prognostic modelling provides a more reliable foundation for future Parkinson's disease stratification research.
  \item \textbf{Open science:} Shared code, synthetic tests, configuration, validation evidence, and figure-generation scripts enable inspection, reuse, and extension.
\end{itemize}
\end{block}

\begin{block}{References}
\nocite{ppmi_dataset,mckinney2010,pedregosa2011}
\footnotesize
\bibliographystyle{plainnat}
\bibliography{poster}
\end{block}

% Deliberate breathing room places the metadata panel level with the
% lower edge of the central results block on the 120 cm x 72 cm canvas.
\vspace{11cm}

\begin{block}{Open Science \& Metadata}
\begin{center}
  \begin{minipage}[t]{0.46\linewidth}
    \centering
    \href{https://github.com/farishaq626-cloud/parkinsons_pipeline}{%
      \includegraphics[width=5.0cm]{github_qr.png}%
    }\\[0.25cm]
    \textbf{GitHub repository}
  \end{minipage}
  \hfill
  \begin{minipage}[t]{0.46\linewidth}
    \centering
    \href{https://doi.org/10.5281/zenodo.21225810}{%
      \includegraphics[width=5.0cm]{zenodo_qr.png}%
    }\\[0.25cm]
    \textbf{Zenodo archive}
  \end{minipage}

  \vspace{0.5cm}

  \textbf{PPMI-Pipeline v2.1.6}\\[0.25cm]
  {\fontsize{15}{19}\selectfont
    \href{https://github.com/farishaq626-cloud/parkinsons_pipeline}{%
      \texttt{github.com/farishaq626-cloud/parkinsons\_pipeline}%
    }\\
    \href{https://doi.org/10.5281/zenodo.21225810}{%
      DOI 10.5281/zenodo.21225810%
    }
  }
\end{center}
\end{block}

\end{column}

\separatorcolumn

\end{columns}
\end{frame}

\end{document}
